
où @@M0@@ est le @@M1@@-schéma
obtenu en faisant éclater @@M2@@. On suppose en outre que @@M3@@ est quasi-compact et a un faisceau ample.

\textbf{Proposition 4.2.} \uline{Le morphisme} @@M4@@ \uline{est projectif}, \uline{d'intersection complète}, \uline{et de dimension relative virtuelle nulle}. \uline{On a de plus la relation} :
@@M5@@

La première assertion découle immédiatement de (VII 1.8).

La deuxième relation est conséquence de la relation :
@@M6@@
montrée dans (VII 3.6).

\textbf{Proposition 4.3.} \uline{Pour tout} @@M7@@, \uline{on a la relation} :
@@M8@@
\uline{où} @@M9@@ \uline{désigne la} @@M10@@-\uline{ième classe de Chern}.

Cette relation est contenue dans VII 4.8, car elle n'est autre que @@M11@@.

\textbf{Proposition 4.4.} \uline{Soit} @@M12@@ \uline{un schéma quasi-compact}. \uline{Si} @@M13@@ \uline{est un faisceau localement libre de type fini sur} @@M14@@, \uline{de classe} @@M15@@ \uline{dans} @@M16@@, \uline{et de rang} @@M17@@, \uline{on a} :
@@M18@@

Les deux membres de la relation à démontrer sont multiplicatifs : c'est clair pour @@M19@@, et pour @@M20@@; c'est aussi vrai pour @@M21@@, qui est la classe de @@M22@@. Par suite, on peut, en passant au fibré en drapeaux de @@M23@@, supposer que @@M24@@ est somme de classes de faisceaux inversibles, et se ramener à montrer la formule pour un faisceau inversible. Soit
